% Paquets généraux
\documentclass[a4paper,12pt,titlepage,twoside]{article}
\usepackage[utf8]{inputenc}
\usepackage[T1]{fontenc}
\usepackage[french]{babel}
\usepackage[gen]{eurosym}
%\usepackage[dvips]{graphicx}
\usepackage{fancyhdr}
\usepackage{pdfpages}
\usepackage{multido}
\usepackage{hyperref}
\usepackage{numprint}
\nplpadding{2}
\usepackage{pstricks}
%\usepackage{aeguill}

\newcommand{\id}{11}
\newcommand{\nom}{Analyse de la structure d'un système asservi}
\newcommand{\sequence}{02}
\newcommand{\num}{02}
\newcommand{\type}{C}
\newcommand{\descrip}{Définition et structure des systèmes asservis, insertion dans les chaînes d'information et d'énergie, les critères de qualité d'un système asservi}
\newcommand{\competences}{A3-C11: Structure des systèmes asservis \\ &  Mod2-C4: Systèmes linéaires continus invariants asservis \\ &  Mod2-C20: Modélisation des systèmes asservis}
\newcommand{\nbcomp}{3}
\newcommand{\systemes}{}
\newcommand{\ilot}{0}

\newcommand\twodigits[1]{%
   \ifnum#1<10 0#1\else #1\fi
}

\newcommand{\auteurun}{Renaud Costadoat}
\newcommand{\auteurdeux}{Françoise Puig}
\newcommand{\auteurs}{Costadoat - Puig}
\newcommand{\institute}{Lycée Dorian}
\newcommand{\scriptpython}{\href{https://raw.githubusercontent.com/Costadoat/S\sequence/master/\type\num\space\nom/Ilot_\twodigits{\ilot} \systemes/\sequence-\type\num-I\twodigits{\ilot}.py}{script python }}

\usepackage{color}
\usepackage{xcolor}
\usepackage{colortbl}
\usepackage{helvet}
\renewcommand{\familydefault}{\sfdefault}
\usepackage{amsfonts}
\usepackage{amsmath}
%\usepackage{xspace}
\usepackage{varioref}
\usepackage{tabularx}
%\usepackage{floatflt}
\usepackage{graphics}
\usepackage{wrapfig}
\usepackage{textcomp}
\usepackage{tikz}
\usepackage{wrapfig}
\usepackage{gensymb}
\usepackage[european]{circuitikz}
\usetikzlibrary{babel}
\usepackage{ifthen}
\usepackage{cancel}
\usepackage{etoolbox}
\usepackage{multirow}
%\usepackage{boxedminipage}
\definecolor{gris25}{gray}{0.75}
\definecolor{analyser}{RGB}{196,189,151}
\definecolor{analyserf}{RGB}{73,69,41}
\definecolor{modeliser}{RGB}{252,213,180}
\definecolor{modeliserf}{RGB}{226,107,10}
\definecolor{resoudre}{RGB}{216,228,188}
\definecolor{resoudref}{RGB}{118,147,60}
\definecolor{experimenter}{RGB}{230,184,183}
\definecolor{experimenterf}{RGB}{150,54,52}
\definecolor{concevoir}{RGB}{141,180,226}
\definecolor{concevoirf}{RGB}{15,36,62}
\definecolor{realiser}{RGB}{184,204,228}
\definecolor{realiserf}{RGB}{54,96,146}
\definecolor{communiquer}{RGB}{204,192,218}
\definecolor{communiquerf}{RGB}{96,73,122	}
\definecolor{bleu}{RGB}{18,33,98}
\definecolor{bleutf}{RGB}{8,53,154}
\definecolor{bleuf}{RGB}{42,94,171}
\definecolor{bleuc}{RGB}{231,239,247}
\definecolor{rougef}{RGB}{185,18,27}
\definecolor{rougec}{RGB}{255,188,204}%255,230,231
\definecolor{vertf}{RGB}{103,126,82}
\definecolor{vertc}{RGB}{220,255,191}
\definecolor{forestgreen}{rgb}{0.13,0.54,0.13}
\definecolor{blcr}{rgb}{0.59,0.69,0.84}
\definecolor{blfr}{rgb}{0.32,0.51,0.75}
\definecolor{orfr}{rgb}{0.90,0.42,0.15}
\definecolor{orcr}{rgb}{0.90,0.65,0.50}
\definecolor{orangef}{rgb}{0.89,0.56,0.16}
\definecolor{orange}{rgb}{0.58,0.35,0.063}
\definecolor{orangec}{rgb}{0.94,0.76,0.56}
\definecolor{rcorrect}{rgb}{0.6,0,0}
\definecolor{sequence}{rgb}{0.75,0.75,0.75}
\definecolor{competences}{rgb}{0.61,0.73,0.35}
\definecolor{grisf}{HTML}{222222}
\definecolor{grisc}{HTML}{636363}
\definecolor{normal}{HTML}{4087c4}
\definecolor{info}{HTML}{5bc0de}
\definecolor{success}{RGB}{92,184,92}
\definecolor{warning}{RGB}{240,173,78}
\definecolor{danger}{RGB}{217,83,79}
\hypersetup{                    % parametrage des hyperliens
    colorlinks=true,                % colorise les liens
    breaklinks=true,                % permet les retours à la ligne pour les liens trop longs
    urlcolor= blfr,                 % couleur des hyperliens
    linkcolor= orange,                % couleur des liens internes aux documents (index, figures, tableaux, equations,...)
    citecolor= forestgreen                % couleur des liens vers les references bibliographiques
    }

% Mise en page
\pagestyle{fancy}

\setlength{\hoffset}{-18pt}

\setlength{\oddsidemargin}{0pt} 	% Marge gauche sur pages impaire2s
\setlength{\evensidemargin}{0pt} 	% Marge gauche sur pages paires
\setlength{\marginparwidth}{00pt} 	% Largeur de note dans la marge
\setlength{\headwidth}{481pt} 	 	% Largeur de la zone de tête (17cm)
\setlength{\textwidth}{481pt} 	 	% Largeu\textbf{r de la zone de texte (17cm)
\setlength{\voffset}{-18pt} 		% Bon pour DOS
\setlength{\marginparsep}{7pt}	 	% Séparation de la marge
\setlength{\topmargin}{-30pt} 		% Pas de marge en haut
\setlength{\headheight}{45pt} 		% Haut de page
\setlength{\headsep}{20pt} 			% Entre le haut de page et le texte
\setlength{\footskip}{30pt} 		% Bas de\textbf{ page + séparation
\setlength{\textheight}{700pt} 		% Hauteur de l'icone zone de texte (25cm)
\setlength\fboxrule{1 pt}
\renewcommand{\baselinestretch}{1}
\setcounter{tocdepth}{1}
\newcommand{\cadre}[2]
{\fbox{
  \begin{minipage}{#1\linewidth}
   \begin{center}
    #2\\
   \end{center}
  \end{minipage}
 }
}

\newcounter{num_quest} \setcounter{num_quest}{0}
\newcounter{num_rep} \setcounter{num_rep}{0}
\newcounter{num_cor} \setcounter{num_cor}{0}

\newcommand{\question}[1]{\refstepcounter{num_quest}\par
~\ \\ \parbox[t][][t]{0.15\linewidth}{\textbf{Question \arabic{num_quest}}}\parbox[t][][t]{0.93\linewidth}{#1}\par
}


\newcommand{\reponse}[1]
{\refstepcounter{num_rep}
\noindent
\rule{\linewidth}{.5pt}
\textbf{Question \arabic{num_rep}:}
\multido{\i=1+1}{#1}{~\ \\}
}

\newcommand{\cor}
{\refstepcounter{num_cor}
\noindent
\rule{\linewidth}{.5pt}
\textbf{Question \arabic{num_cor}:} \\
}

\newcommand{\titre}[1]
{\begin{center}
\cadre{0.8}{\huge #1}
\end{center}
}

%Titres des sections

\newcommand{\titleana}[2]{%
\tikzstyle{titlebox}=[rectangle,inner sep=10pt,inner ysep=10pt,draw,fill=analyser]%
\tikzstyle{title}=[fill=analyser]%
%
\bigskip\noindent\begin{tikzpicture}
\node[titlebox] (box){%
\begin{minipage}{0.94\textwidth}
\textbf{\textcolor{analyserf}{#1}}
\end{minipage}
};
%\draw[color=blue] (box.north west)--(box.north east);
\node[title] at (box.north) {\textcolor{analyserf}{ANALYSER}};
\end{tikzpicture}%
\\}

\newcommand{\titlemod}[2]{%
\tikzstyle{titlebox}=[rectangle,inner sep=10pt,inner ysep=10pt,draw,fill=modeliser]%
\tikzstyle{title}=[fill=modeliser]%
%
\bigskip\noindent\begin{tikzpicture}
\node[titlebox] (box){%
\begin{minipage}{0.94\textwidth}
\textbf{\textcolor{modeliserf}{#1}}
\end{minipage}
};
%\draw[color=blue] (box.north west)--(box.north east);
\node[title] at (box.north) {\textcolor{modeliserf}{MODELISER}};
\end{tikzpicture}%
\\}

\newcommand{\titleres}[2]{%
\tikzstyle{titlebox}=[rectangle,inner sep=10pt,inner ysep=10pt,draw,fill=resoudre]%
\tikzstyle{title}=[fill=resoudre]%
%
\bigskip\noindent\begin{tikzpicture}
\node[titlebox] (box){%
\begin{minipage}{0.94\textwidth}
\textbf{\textcolor{resoudref}{#1}}
\end{minipage}
};
%\draw[color=blue] (box.north west)--(box.north east);
\node[title] at (box.north) {\textcolor{resoudref}{RESOUDRE}};
\end{tikzpicture}%
\\}

\newcommand{\titleexp}[2]{%
\tikzstyle{titlebox}=[rectangle,inner sep=10pt,inner ysep=10pt,draw,fill=experimenter]%
\tikzstyle{title}=[fill=experimenter]%
%
\bigskip\noindent\begin{tikzpicture}
\node[titlebox] (box){%
\begin{minipage}{0.94\textwidth}
\textbf{\textcolor{experimenterf}{#1}}
\end{minipage}
};
%\draw[color=blue] (box.north west)--(box.north east);
\node[title] at (box.north) {\textcolor{experimenterf}{EXPERIMENTER}};
\end{tikzpicture}%
\\}

\newcommand{\titlecon}[2]{%
\tikzstyle{titlebox}=[rectangle,inner sep=10pt,inner ysep=10pt,draw,fill=concevoir]%
\tikzstyle{title}=[fill=concevoir]%
%
\bigskip\noindent\begin{tikzpicture}
\node[titlebox] (box){%
\begin{minipage}{0.94\textwidth}
\textbf{\textcolor{concevoirf}{#1}}
\end{minipage}
};
%\draw[color=blue] (box.north west)--(box.north east);
\node[title] at (box.north) {\textcolor{concevoirf}{CONCEVOIR}};
\end{tikzpicture}%
\\}

\newcommand{\titlerea}[2]{%
\tikzstyle{titlebox}=[rectangle,inner sep=10pt,inner ysep=10pt,draw,fill=realiser]%
\tikzstyle{title}=[fill=realiser]%
%
\bigskip\noindent\begin{tikzpicture}
\node[titlebox] (box){%
\begin{minipage}{0.94\textwidth}
\textbf{\textcolor{realiserf}{#1}}
\end{minipage}
};
%\draw[color=blue] (box.north west)--(box.north east);
\node[title] at (box.north) {\textcolor{realiserf}{REALISER}};
\end{tikzpicture}%
\\}

\newcommand{\titlecom}[2]{%
\tikzstyle{titlebox}=[rectangle,inner sep=10pt,inner ysep=10pt,draw,fill=communiquer]%
\tikzstyle{title}=[fill=communiquer]%
%
\bigskip\noindent\begin{tikzpicture}
\node[titlebox] (box){%
\begin{minipage}{0.94\textwidth}
\textbf{\textcolor{communiquerf}{#1}}
\end{minipage}
};
%\draw[color=blue] (box.north west)--(box.north east);
\node[title] at (box.north) {\textcolor{communiquerf}{COMMUNIQUER}};
\end{tikzpicture}%
\\}


\newcommand{\prob}[2]{%
\tikzstyle{titlebox}=[rounded corners=10,inner sep=10pt,inner ysep=10pt,draw,fill=bleuf]%
%\tikzstyle{title}=[]%
%
\bigskip\noindent\begin{tikzpicture}
\node[titlebox] (box){%
\begin{minipage}{0.94\textwidth}
\begin{minipage}{0.15\linewidth}
\includegraphics[width=0.8\linewidth]{../../../img/questionmark}
\end{minipage}
\begin{minipage}{0.8\linewidth}
\textbf{\textcolor{white}{
 Problématique du TP:\\ \noindent\rule{13cm}{1pt} \\ \\#1}}
\end{minipage}
\end{minipage}
};
%\draw[color=blue] (box.north west)--(box.north east);
%\node[title] at (box.north west) {\textcolor{white}{\textbf{Problématique}}};
\end{tikzpicture}\bigskip\\~\ \\%
}


% En tête et pied de page
\lhead{\nom}
\rhead{\includegraphics[width=2cm]{../../../img/logo}}
\lfoot{\auteurs}
\cfoot{Page \thepage}

\fancypagestyle{correction}{%
  \fancyhf{}
  \lhead{\colorbox{danger}{\begin{minipage}{0.65\paperwidth} \textcolor{white}{\textbf{Correction}} \end{minipage}} }
  \rhead{\includegraphics[width=2cm]{../../../img/logo}}
  \lfoot{Renaud Costadoat}
  \rfoot{\colorbox{danger}{\begin{minipage}{0.6\paperwidth} \begin{flushright}\textcolor{white}{\textbf{Correction}}\end{flushright} \end{minipage}} }}

\renewcommand{\footrulewidth}{0.4pt}

\usepackage{eso-pic}
\newcommand{\BackgroundPic}{%
\put(0,0){%
\parbox[b][\paperheight]{\paperwidth}{%
\vfill
\begin{center}
\hspace{0.5cm}\vspace{0.5cm}
\includegraphics[width=\paperwidth,height=\paperheight,%
keepaspectratio]{../../../img/fond3}%
\end{center}
\vfill
}}}

\newcommand{\BackgroundPicdeux}{%
\put(25,-30){%
\parbox[b][\paperheight]{\paperwidth}{%
\vfill
\begin{center}
\includegraphics[width=\paperwidth,height=\paperheight,%
keepaspectratio]{../../../img/fond4}%
\end{center}
\vfill
}}}

\newcommand{\activite}[1]{\newpage \cleardoublepage
\renewcommand\arraystretch{2.5}\setcounter{section}{0}\setcounter{subsection}{0}
\ifthenelse{\equal{#1}{1}}{\begin{center}\begin{tabular}{|c|cc p{3cm}|}
\hline
\rowcolor{violt} \supportu & Act & Contenu & Compétences \\\hline
 \multirow{3}{*}{\includegraphics[width=0.2\linewidth]{\supportimageu}} & \begin{LARGE}\textbf{1}\end{LARGE} & \activiteun & \competenceun \\
 & \multicolumn{3}{l|}{Nom: .........................................................................}\\
 & \multicolumn{3}{l|}{Prénom: ....................................................................}\\\hline
\end{tabular}\end{center}}{}
\ifthenelse{\equal{#1}{2}}{\begin{center}\begin{tabular}{|c|cc p{3cm}|}
\hline
\end{tabular}\end{center}}{}}

\begin{document}

\pagestyle{empty}

\vspace*{-3\baselineskip}

\AddToShipoutPicture*{\BackgroundPic}

\ifdef{\auteurdeux}{\begin{tabular}{>{\columncolor{gray!00}}m{.33\linewidth} m{.27\linewidth} >{\columncolor{gray!00}}m{.3\linewidth}}
Séquence \sequence \ - \type \num \ - Îlot \numprint{\ilot} &  \multirow{3}{*}{\hspace{1cm}\includegraphics[height=1.5cm]{../../../img/logo}} &  \begin{flushright} \multirow{4}{*}{\hspace{1cm}\includegraphics[height=4cm]{img/qrcode}}\end{flushright}\\
 \textbf{\institute} \\
 \auteurun\\
 \auteurdeux
\end{tabular}}{\begin{tabular}{>{\columncolor{gray!00}}m{.3\linewidth} m{.3\linewidth} >{\columncolor{gray!00}}m{.3\linewidth}}
Séquence \sequence \ - \type \num \ - Îlot \numprint{\ilot}  &  \multirow{3}{*}{\hspace{1cm}\includegraphics[height=1.5cm]{../../../img/logo}} &  \begin{flushright} \multirow{4}{*}{\hspace{1cm}\includegraphics[height=4cm]{img/qrcode}}\end{flushright}\\
 \textbf{\institute} \\
 \auteurun
\end{tabular}}

\vspace{1cm}

\begin{center}\colorbox{white}{\huge{\nom}}\end{center}

\vspace{2cm}


\begin{center}\colorbox{white!20}{\includegraphics[height=5cm]{/home/renaud/Documents/Renaud/GitHub/django_education/django_education/static/systemes/\imageun}}\end{center}

\vspace{2cm}



\begin{tabular}{p{.15\linewidth} >{\columncolor{white}}p{.8\linewidth}}
    \rowcolor{gray!20}
    Référence & S\sequence \ - \type\num \ - I\numprint{\ilot} \\
    Compétences & \competences \\
 	\rowcolor{gray!20}
    Description & \descrip \\
    Système & \systemes
  \end{tabular}

\newpage

\AddToShipoutPicture{\BackgroundPicdeux}

\pagestyle{fancy}


\ifdef{\public}{\cleardoublepage}{}

\ifdef{\public}{\cleardoublepage}{}

\prob{Modéliser un Système Linéaire Continu et Invariant} 

\titleexp{Effectuer la mesure du déplacement du charriot} 

La procédure suivante permet de mettre en \oe uvre et de mesurer les conséquences de l'asservissement de l'axe Emericc.

Plusieurs capteurs sont reliés à la carte du boîtier qui permet l'acquisition, le traitement et affichage sont effectués par l'ordinateur. Pour cela :
\begin{enumerate}
 \item Lancer le logiciel winaxe32,
 \item Lancer une recherche de la carte,
 \item Effectuer une prise d'origine grâce à la fonction \og prise d'origine\fg,
 \item Sélectionner le pilotage en \og Boucle ouverte \fg,
 \item Choisir des paramètres de configuration (faire plusieurs expériences en faisant varier les paramètres),
 \item Cliquer sur lancer la mesure.
\end{enumerate} 

\paragraph{Question 1:} Représenter la courbe obtenue sur un graphe sur lequel vous ferez apparaitre la consigne échelon, l'écart statique et le temps de réponse à 5\%.

~\

\textit{Effectuer une nouvelle mesure en ajoutant des masses supplémentaires sur le charriot en translation.}

\paragraph{Question 2:} Représenter la nouvelle courbe obtenue sur un graphe sur lequel vous ferez apparaitre la consigne échelon, l'écart statique et le temps de réponse à 5\%.

\paragraph{Question 3:} En utilisant l'ensemble des mesures mis à votre disposition, expliquer les différences entre les deux courbes.

~\

L'ensemble des réponses que vous aurez donné dans cette partie devra être utilisé afin de compléter le document de présentation. Les diapositives pourront être agrémentées comme vous le souhaitez.

\ifdef{\public}{\cleardoublepage}{\newpage}

\titlemod{Modélisation de la chaine d'énergie du système}

La mise en mouvement du charriot est assurée par le mouvement du moteur. Nous allons dans cette partie chercher à déterminer la fonction de transfert qui lie la tension en entrée du moteur $u(t)$ et le déplacement du charriot $x(t)$. L'inertie du système sera négligée, comme la masse du charriot devant la masse des disques embarqués.

~\

Caractéristiques du moteur:
\begin{tabular}{|c|c|}
\hline
$K_m$ & $34,2.10^{-3}N.m.A^{-1}$ \\
\hline
$R_m$ & $4\Omega$ \\
\hline
$L_m$ & $0,63.10^{-3}H$ \\
\hline
\end{tabular}

~\

\paragraph{Question 4:} Donner l'ensemble des équations temporelles permettant de modéliser un moteur à courant continu (on ne prendra pas en compte son inductance interne).

\paragraph{Question 5:} Déterminer l'équation qui lie la vitesse de rotation du moteur $\omega_m(t)$ à la vitesse de déplacement du charriot $v(t)$. Puis la position du charriot $x(t)$ en fonction de sa vitesse de déplacement $v(t)$. Effectuer sur le système toutes les mesures nécessaires à l'écriture de ces équations.

\paragraph{Question 6:} Passer ces équations dans le domaine de Laplace.

\paragraph{Question 7:} Mettre le système sous la forme de la fonction de transfert suivante: $H(p)=\dfrac{X(p)}{U_m(p)}$. Donner les caractéristiques de cette fonction de transfert et vérifier l'homogénéité des constantes déterminées.

~\

L'ensemble des réponses que vous aurez donné dans cette partie devra être utilisé afin de compléter le document de présentation. Les diapositives pourront être agrémentées comme vous le souhaitez.

\ifdef{\public}{\cleardoublepage}{\newpage}

\titleres{Activité 3: Simulation du comportement du modèle}

Le logiciel \textbf{Scilab} permet de tracer la réponse temporelle d'une fonction de transfert donnée.

Pour cela, il suffit de lancer le logiciel et d'aller dans le module \textbf{Xcos}.

Dans le dossier \textbf{CPGE} du navigateur de palettes, vous trouverez, par exemple:
\begin{itemize}
 \item une \textit{entrée}: STEP\_FUNCTION,
 \item un \textit{Opérateur linéaire}: CLR, vous modifierez sa fonction de transfert afin d'obtenir ce que vous souhaitez observer,
 \item une \textit{sortie}: SCOPE,
 \item un \textit{outil d'analyse}: REP\_TEMP.
\end{itemize}

Faire glisser ces blocs sur une page vierge du module xcos et cliquer sur la flèche permettant de lancer la simulation.

\paragraph{Question 8:} Effectuer le tracé de la fonction de transfert vue en TD afin d'apprendre à maitriser l'outil Scilab.

\paragraph{Question 9:} Modifier la fonction de transfert afin d'y insérer celle modélisée dans l'activité 2.

\paragraph{Question 10:} Critiquer les résultats obtenus et analyser le lien entre les tracés obtenus par simulation et ceux obtenus durant l'activité expérimentale 1.

~\

L'ensemble des réponses que vous aurez donné dans cette partie devra être utilisé afin de compléter le document de présentation. Les diapositives pourront être agrémentées comme vous le souhaitez.

\ifdef{\public}{\end{document}}{}

\newpage

\pagestyle{correction}\setcounter{section}{0}

Modélisation
\begin{center}
$H(p)=\dfrac{X_e(p)}{U_m(p)}=\dfrac{\dfrac{R_p.r}{K_e}}{p.\left(1+\dfrac{R.J}{K_m.K_e}.p+\dfrac{L.J}{K_m.K_e}.p^2\right)}$
\end{center}

\end{document}